% Standalone manuscript. Compile with XeLaTeX or Tectonic; no external assets.
\documentclass[11pt,a4paper]{article}
\usepackage[margin=24mm,headheight=15pt]{geometry}
\usepackage{iftex}
\ifPDFTeX
  \usepackage[T1]{fontenc}\usepackage[utf8]{inputenc}\usepackage{lmodern}
\else
  \usepackage{fontspec}
  \IfFontExistsTF{CMU Serif}{\setmainfont{CMU Serif}}{\setmainfont{Latin Modern Roman}}
  \setsansfont{NotoSans}[Extension=.ttf,UprightFont=*-Regular,BoldFont=*-Bold,
    ItalicFont=*-Italic,BoldItalicFont=*-BoldItalic,Scale=MatchLowercase]
\fi
\usepackage{amsmath,amssymb,amsthm,mathtools,microtype}
\usepackage{booktabs,tabularx,array,enumitem,xcolor,titlesec,fancyhdr,needspace}
\usepackage{tikz}
\usetikzlibrary{arrows.meta,positioning}
\definecolor{LinkGreen}{HTML}{176B4A}
\definecolor{LinkPurple}{HTML}{70459B}
\definecolor{Ink}{HTML}{20384D}
\usepackage[colorlinks=true,linkcolor=LinkGreen,citecolor=LinkGreen,urlcolor=LinkPurple,
  bookmarksnumbered=true,pdfusetitle]{hyperref}
\usepackage{xurl}
\urlstyle{same}
\newtheorem{theorem}{Theorem}[section]
\newtheorem{lemma}[theorem]{Lemma}
\newtheorem{proposition}[theorem]{Proposition}
\newtheorem{corollary}[theorem]{Corollary}
\theoremstyle{definition}\newtheorem{definition}[theorem]{Definition}
\newtheorem{example}[theorem]{Example}
\theoremstyle{remark}\newtheorem{remark}[theorem]{Remark}
\newcommand{\E}{\mathbb E}\newcommand{\Prob}{\mathbb P}
\newcommand{\R}{\mathbb R}\newcommand{\C}{\mathbb C}
\newcommand{\Tr}{\operatorname{Tr}}\newcommand{\tr}{\operatorname{Tr}}
\newcommand{\norm}[1]{\lVert#1\rVert}\newcommand{\ip}[2]{\langle#1,#2\rangle}
\newcommand{\diag}{\operatorname{diag}}
\numberwithin{equation}{section}
\titleformat{\section}{\large\sffamily\bfseries}{\thesection}{.6em}{}
\titleformat{\subsection}{\normalsize\sffamily\bfseries}{\thesubsection}{.6em}{}
\titlespacing*{\section}{0pt}{15pt}{7pt}
\titlespacing*{\subsection}{0pt}{10pt}{5pt}
\setlength{\parindent}{1.1em}\setlength{\parskip}{3pt}
\setlength{\emergencystretch}{2em}
\setlist{itemsep=3pt,topsep=4pt}
\allowdisplaybreaks[1]\raggedbottom
\clubpenalty=10000\widowpenalty=10000
\pagestyle{fancy}\fancyhf{}
\fancyhead[L]{\sffamily\footnotesize Finite-moment sketch refinements}
\fancyhead[R]{\sffamily\footnotesize SRHT and Khatri--Rao}
\fancyfoot[C]{\sffamily\small\thepage}
\renewcommand{\headrulewidth}{.3pt}
\setcounter{tocdepth}{1}
\title{Finite-moment refinements for\\SRHT and Khatri--Rao sketches}
\author{Diar Heidary}
\date{5 October 2026}
\begin{document}
\maketitle
\begin{abstract}
We refine finite-moment bounds for two structured subspace embeddings by
retaining information that is lost in a terminal operator-norm estimate.
For two-round SRHT, successive polynomial prefixes use smaller sign
cutoffs, yielding an ordered product of nonnegative comparison matrices.
For Gaussian Khatri--Rao sketches, an ordered first-return comparison
retains the covariance spectrum at the vacuum boundary, and exact
fourth-order calibration gives additional finite-sample certificates.
We also transfer every even trace-moment bound for Gaussian product
probes to independent normalized spherical product probes by conditional
Jensen, compute their exact covariance reduction, and give an exact
fourth-moment example. A conditional spherical Jacobi comparison explains
the related linear-pencil case. The proofs identify the estimates
imported from earlier manuscripts and distinguish finite-parameter gains
from changes to universal asymptotic constants.
\end{abstract}
\tableofcontents
\section{Scope and contributions}\label{sec:introduction}
For a self-adjoint random matrix $Z$, an even trace moment gives the
sufficient tail certificate
\begin{equation}\label{eq:markov}
 \Prob\{\norm Z>\varepsilon\}
 \le \varepsilon^{-2q}\E\Tr Z^{2q},\qquad q\ge1.
\end{equation}
Throughout, $\Tr$ is the unnormalized trace, $\norm\cdot$ is the operator
norm, and comparison vectors use the Euclidean norm. A sufficient sample
count obtained from \eqref{eq:markov} is not asserted to be an optimal
threshold.

The common starting point is multiplication by the original random
matrix on vector-valued $L^2$. The constant functions form the vacuum.
A cutoff containing the first $q$ vacuum iterates preserves the
$2q$-th trace moment. The estimates below keep the probability law of the
original SRHT or Gaussian Khatri--Rao sketch; the spherical-product
section explicitly introduces a different probe law.

\begin{center}
\small
\begin{tabularx}{\linewidth}{@{}>{\raggedright\arraybackslash}p{.22\linewidth}>{\raggedright\arraybackslash}X >{\raggedright\arraybackslash}X@{}}
\toprule
Setting & Information retained & Resulting refinement\\
\midrule
Two-round SRHT & Moment-step cutoffs and selector occupation & Ordered-prefix moment and row certificates\\
Gaussian Khatri--Rao & Covariance eigenvalues and exact fourth-order reuse & Boundary comparison and calibrated fourth moments\\
Spherical products & Gaussian radius--direction decomposition & All-order transfer and a sharper covariance\\
Spherical linear pencils & Conditional covariance and Hermite grade costs & Finite Jacobi comparison\\
\bottomrule
\end{tabularx}
\end{center}

\paragraph{Provenance.}
The original-law multiplication representation, positive-selector
estimates, and ordered covariance-spectrum comparison are developed in
\cite{GradedSource}. The SRHT row-factor estimates come from
\cite{SRHTSource}, and the nonlinear Khatri--Rao repeated-visit envelope
comes from \cite{KRSource}. The Gaussian creation estimate used here is
due to Bandeira, Kunisky, Nizi\'{c}-Nikolac, Pesenti, and Wang
\cite{Bandeira}. Gaussian Wick contraction, Gaussian polar decomposition,
conditional Jensen, and the Schur-complement renewal identity are
classical ingredients. Our purpose is to prove the stated refinements
and extensions from these inputs; no exhaustive literature-priority
claim is made for the combinations. These new statements have not been
formally verified in Lean.

\paragraph{What the numerical comparisons establish.}
Every displayed row count is tied to a specified moment order, model,
and comparison bound. The SRHT table improves a terminal-grade
certificate. The covariance-spectrum Khatri--Rao table improves a
stationary scalar-boundary certificate, while the fourth-order table
compares fourth-order certificates. None of these comparisons alone
replaces a universal leading constant in the strongest available
all-order theorem.


% Required macros: \E, \Tr, \norm{...}; theorem/lemma/corollary/remark environments.
% Bibliography keys: SRHTSource, GradedSource.
\section{Two-round SRHT: selector returns and ordered prefixes}
\label{sec:srht-prefix}

The refinement in this section changes how the existing local row estimates
are assembled. It retains selector grades, then uses the smaller sign spaces
reached by successive prefixes of the moment word. The distribution of the
sketch, its coherence parameters, and the local coefficient contractions are
unchanged. In particular, the argument uses the original sign law.

Let $A,B\in\mathbb C^{n\times n}$ be unitary and let
$U\in\mathbb C^{n\times d}$ satisfy $U^*U=I_d$. Define
\[
 \alpha=n\max_{r,a}|A_{ra}|^2,\qquad
 \beta=n\max_{a,i}|B_{ai}|^2,\qquad
 W=AD_yBD_xU,
\]
where $D_x,D_y$ are independent diagonal matrices of independent unbiased
signs. Write $w_r^*=e_r^*W$. Independently choose a uniform $m$-element
subset $T\subseteq[n]$, and set
\begin{equation}
 E_T=\frac nm\sum_{r\in T}w_rw_r^*-I_d.
 \label{eq:srht-error}
\end{equation}
For flat unitary transforms, including normalized Walsh transforms,
$\alpha=\beta=1$.

\subsection{The finite-moment certificate}

Fix an integer $q\ge1$ and $1\le m<n$, and write
\[
 \rho=m/n,\qquad u=\sqrt{(1-\rho)/\rho},\qquad
 z=|1-2\rho|/\rho,
 \qquad
 \underline c=\max\left\{\frac12,
       1-\frac1{2\sqrt{m(1-m/n)}}\right\}.
\]
For $1\le j\le q$ and $1\le k\le\min(j,n)$, define
\begin{align}
 \Lambda_{j,k}
 &=\left(\sqrt{\alpha k(2j+1)}
      +\sqrt{2\{\beta(d+8j^2)+\alpha k(2j-1)\}}\right)^2,
       \label{eq:srht-Lambda}\\
 \ell_{j,k}&=\min\{1,\Lambda_{j,k}/n\},\qquad \ell_{j,0}=0.
       \label{eq:srht-ell}
\end{align}
Let $G^{[j]}$ be the symmetric nonnegative matrix on indices $0,\ldots,q$
whose nonzero entries are confined to $0,\ldots,\min(j,n)$ and satisfy
\begin{equation}
 G^{[j]}_{kk}=z\ell_{j,k},\qquad
 G^{[j]}_{k,k+1}=G^{[j]}_{k+1,k}=u\sqrt{\ell_{j,k+1}}.
 \label{eq:srht-stage-matrix}
\end{equation}
Here and below $e_0$ is the first standard basis vector of this scalar
comparison space.

\begin{theorem}[Ordered-prefix SRHT certificate]
\label{thm:srht-prefix}
Under the preceding assumptions,
\begin{equation}
 \boxed{\displaystyle
 \E\Tr E_T^{2q}
 \le d\,\underline c^{-2q}
       \bigl\|G^{[q]}G^{[q-1]}\cdots G^{[1]}e_0\bigr\|_2^2.}
 \label{eq:srht-prefix-moment}
\end{equation}
Consequently the right-hand side divided by $\varepsilon^{2q}$ bounds
$\mathbb P\{\norm{E_T}>\varepsilon\}$. Moreover,
\begin{equation}
 \bigl\|G^{[q]}\cdots G^{[1]}e_0\bigr\|_2^2
 \le\bigl\|(G^{[q]})^qe_0\bigr\|_2^2
 =\bigl((G^{[q]})^{2q}\bigr)_{00}.
 \label{eq:srht-terminal-comparison}
\end{equation}
Thus the ordered-prefix certificate is never worse than the terminal-cutoff
selector certificate with the same parameters. For $m=n$, $E_T=0$ exactly.
The conclusions also hold when each sign family is
$\min(n,4q)$-wise independent, provided the two families remain mutually
independent and independent of $T$.
\end{theorem}

The order in \eqref{eq:srht-prefix-moment} matters. The matrices
$G^{[j]}$ need not commute, and the squared norm of their product is not
generally a vacuum entry of the power of one matrix. The prefactor remains
$d$: the retained sign space does not supply extra initial trace states.

\subsection{Local row estimates and the legal cutoff}

We record explicitly which local estimates are inherited from the existing
proof. In the Boolean character basis of the two sign families, multiplication
by $x_i$ is $X_i=a_{x,i}+a_{x,i}^*$, where creation inserts an unoccupied
index and annihilation removes an occupied index. In particular,
$X_i^*X_i=I$. Sign grades count occupied coordinates.

Let $P_{\mathrm{sig},j}$ retain grades at most $\min(2j,n)$ in each sign
family, with the external $\mathbb C^d$ factor included. On its range define
the positive effects
\begin{equation}
 A_r^{[j]}=P_{\mathrm{sig},j}\,M_{w_rw_r^*}\,P_{\mathrm{sig},j},
 \qquad \sum_{r=1}^n A_r^{[j]}=I,
 \label{eq:srht-effects}
\end{equation}
where $M_F$ denotes pointwise multiplication by the matrix-valued function
$F$. The last identity follows from $W^*W=I_d$ pointwise.

\begin{lemma}[Subset allowance at a prefix cutoff]
\label{lem:srht-prefix-subsets}
For $S\subseteq[n]$ with $|S|\le k\le\min(j,n)$,
\begin{equation}
 \left\|\sum_{r\in S}A_r^{[j]}\right\|\le\ell_{j,k}.
 \label{eq:srht-prefix-subsets}
\end{equation}
\end{lemma}
\begin{proof}
Let $\mathcal R_S$ be the row multiplication map
$(\mathcal R_Sf)_r=w_r^*f$, for $r\in S$, with the row label retained as
an orthogonal output coordinate. Crucially, its output sign grades are not
truncated. Then
\[
 \sum_{r\in S}A_r^{[j]}
   =(\mathcal R_SP_{\mathrm{sig},j})^*
     (\mathcal R_SP_{\mathrm{sig},j}).
\]
Decompose $\mathcal R_S=\mathcal C_S+\mathcal A_S$ according to creation
or annihilation in the second sign family, keeping the two first-family
transitions joined in $\mathcal C_S$.

Here are the precise coefficient estimates used. Set $c_i=e_i^*U$ and
$(g_{ri})_a=\sqrt n\,\overline{A_{ra}B_{ai}}$. Unitarity gives
\[
 \left\|\sum_i t_i g_{ri}\right\|^2\le\alpha\sum_i|t_i|^2,
 \qquad
 \sum_r|g_{ri}\rangle\langle g_{ri}|\preceq\beta I,
 \qquad \sum_i c_i^*c_i=I_d.
\]
These inequalities hold with passive Hilbert-space factors. Before a
second-family creation, the first inequality and $X_i^*X_i=I$ bound the
whole mixed-first-grade coefficient map by $\sqrt\alpha$ for each row.
Insertion into second-family grade $s$ costs $\sqrt{s+1}$. Orthogonality
of the output grades therefore yields
\begin{equation}
 n\norm{\mathcal C_SP_{\mathrm{sig},j}}^2
       \le\alpha k(2j+1).
 \label{eq:srht-creation-row}
\end{equation}

For completeness, the fixed-input annihilation estimates inherited from
the coefficient-map and Boolean-kernel arguments of
\cite{GradedSource,SRHTSource} are, on input grades $(r,s)$,
\begin{align*}
 n\norm{\mathcal R_S^{+-}}^2&\le\beta(r+1)s,\\
 n\norm{\mathcal R_S^{--}}^2
      &\le\beta(d+p+pv)+\alpha |S|v,
      \qquad p=(r-1)_+,\quad v=(s-1)_+.
\end{align*}
The superscripts are the signs of the changes in the two sign grades.
The first estimate factors through a second-family extraction and a
first-family insertion. The second is the collective double-annihilation
estimate, including the original Boolean occupation correction; this
technical kernel estimate is an input from the cited proofs, not a Gaussian
replacement or a new assertion here.

For fixed input grades these two outputs are orthogonal. Each fixed output
grade under $\mathcal A_S$ has at most two possible input grades, so
Cauchy--Schwarz contributes a factor of two when mixed input grades are
assembled. Since $r,s\le2j$ implies
$p+pv+(r+1)s\le8j^2$, this gives
\begin{equation}
 n\norm{\mathcal A_SP_{\mathrm{sig},j}}^2
    \le2\{\beta(d+8j^2)+\alpha k(2j-1)\}.
 \label{eq:srht-annihilation-row}
\end{equation}
The triangle inequality applied after
\eqref{eq:srht-creation-row}--\eqref{eq:srht-annihilation-row} proves
the allowance $\Lambda_{j,k}/n$. Positivity and
\eqref{eq:srht-effects} cap it at one.
\end{proof}

This is the existing retained row-factor argument with two parameters
kept separate: $j$ controls the sign cutoff, while $k$ controls the number
of selected row effects. Substituting $k$ for $j$ throughout
\eqref{eq:srht-Lambda} is unjustified: a low selector grade can coexist
with a high sign grade.

\subsection{Proof of the moment and sampling comparisons}

\begin{proof}[Proof of Theorem~\ref{thm:srht-prefix}]
First use independent Bernoulli selectors $\eta_r$ of mean $\rho$,
independent of the signs, and put
\[
 E_\rho=\rho^{-1}\sum_r\eta_rw_rw_r^*-I_d.
\]
In the orthonormal basis
$1,(\eta_r-\rho)/\sqrt{\rho(1-\rho)}$, multiplication by
$\eta_r/\rho-1$ is
\[
 \begin{pmatrix}0&u\\u&\theta\end{pmatrix},
 \qquad \theta=(1-2\rho)/\rho.
\]
Selector grades are therefore subset cardinalities. On the stage-$j$
sign space and selector grades at most $\min(j,n)$, the compressed
multiplier has the exact form
\[
 M_j=u(K_j+K_j^*)+\theta D_j,
\]
where
\[
 (K_jf)_S=\sum_{r\in S}A_r^{[j]}f_{S\setminus\{r\}},
 \qquad
 (D_jf)_S=\left(\sum_{r\in S}A_r^{[j]}\right)f_S.
\]

We give the positive-effects estimate behind its block bounds. For an
output set $S$ of size $k+1$,
\[
 \left\|\sum_{r\in S}A_r^{[j]}v_r\right\|^2
 \le\left\|\sum_{r\in S}A_r^{[j]}\right\|
       \sum_{r\in S}\langle v_r,A_r^{[j]}v_r\rangle.
\]
This follows by factoring the row through the square roots of the
positive effects. Sum over $S$, noting that each input set $J$ is charged
$\sum_{r\notin J}A_r^{[j]}\preceq I$. Lemma~\ref{lem:srht-prefix-subsets}
then gives
\[
 \norm{K_j:\text{grade }k\to k+1}\le\sqrt{\ell_{j,k+1}},
 \qquad
 \norm{D_j:\text{grade }k\to k}\le\ell_{j,k}.
\]
The reverse edge has the same bound by adjointness. These are precisely
the entries of \eqref{eq:srht-stage-matrix} after taking absolute values.

Let $M$ denote the uncompressed multiplication operator for $E_\rho$,
and let $\Omega_a$ be the constant vector $e_a$, with both sign families
and the selector family in their vacuum states. Multiplication raises
each sign grade by at most two and the selector grade by at most one.
Hence $f_j=M^j\Omega_a$ lies in the stage-$j$ space, and, with the
nested spaces identified isometrically,
\[
 f_j=M_j f_{j-1}.
\]
This identity is asserted for these actual prefixes; it does not assert
that compressions preserve arbitrary products. Let $b_j(k)$ be the norm
of the selector-grade-$k$ component of $f_j$, including all its sign
coordinates. The block bounds imply componentwise
\[
 b_j\le G^{[j]}b_{j-1},\qquad b_0=e_0.
\]
All entries are nonnegative, so iteration and orthogonality give
\begin{equation}
 \E\Tr E_\rho^{2q}
   =\sum_{a=1}^d\norm{M^q\Omega_a}^2
   \le d\,\bigl\|G^{[q]}\cdots G^{[1]}e_0\bigr\|_2^2.
 \label{eq:srht-bernoulli-prefix}
\end{equation}

We now transfer to the original fixed-size law, spelling out the centered
coupling because it is distinct from a with-replacement comparison. Draw
a Bernoulli subset $\mathcal B$ of size $K\sim\mathrm{Bin}(n,\rho)$.
If $K>m$, remove uniformly chosen elements until $m$ remain; if $K<m$,
add uniformly chosen missing elements. The resulting $T$ is a uniform
$m$-subset and is independent of $K$. Set
\[
 a=\E(K-m)_+=\E(m-K)_+,\qquad v=m(1-m/n).
\]
The conditional inclusion probabilities of $\mathcal B$ inside and
outside $T$ are $1-a/m$ and $a/(n-m)$, respectively. Consequently, for
arbitrary deterministic Hermitian matrices $H_r$,
\[
 Z_{\mathcal B}=\sum_r(\mathbf1_{r\in\mathcal B}-\rho)H_r,
 \quad Z_T=\sum_r(\mathbf1_{r\in T}-\rho)H_r,
 \qquad
 \E[Z_{\mathcal B}\mid T]=cZ_T,
 \quad c=1-a/v.
\]
The binomial identity $k\binom nk=n\binom{n-1}{k-1}$ shows that
\[
 a=v\,\mathbb P\{\mathrm{Bin}(n-1,\rho)=m\}.
\]
The atoms at $m-1$ and $m$ of this latter binomial law are equal, so
$c\ge1/2$. Also
$a=\frac12\E|K-m|\le\sqrt v/2$, proving $c\ge\underline c$.
Convexity of $Z\mapsto\Tr Z^{2q}=\norm{Z}_{S_{2q}}^{2q}$ on Hermitian
matrices and conditional Jensen yield
\[
 \E\Tr Z_T^{2q}\le c^{-2q}\E\Tr Z_{\mathcal B}^{2q}.
\]
Apply this conditionally on the signs with $H_r=w_rw_r^*$, then rescale
by $\rho^{-1}$ and use \eqref{eq:srht-bernoulli-prefix}. This proves
\eqref{eq:srht-prefix-moment}.

For fixed $k\ge1$, $\Lambda_{j,k}$ is nondecreasing in $j$. After zero
embedding, $0\le G^{[j]}\le G^{[q]}$ entrywise for every $j\le q$.
Nonnegative induction proves \eqref{eq:srht-terminal-comparison}.
Finally, $\Tr E_T^{2q}$ has degree at most $4q$ in each sign family.
Moment matching therefore transfers the fully independent proof to the
stated limited-independent signs. Markov's inequality and
$\norm{E_T}^{2q}\le\Tr E_T^{2q}$ finish the proof.
\end{proof}

\subsection{Certified sufficient sample counts}

Table~\ref{tab:srht-prefix-counts} keeps the same moment order, sampling
transfer, and target failure probability in both columns. It compares
the two sides of \eqref{eq:srht-terminal-comparison}, not against a
coarser universal scalar prescription.

\begin{table}[htbp]
\centering\small
\begin{tabular}{rrrrr}
\toprule
$d$ & $q$ & Terminal cutoff & Ordered prefixes & Reduction\\
\midrule
$10$ & $6$ & $28\,080$ & $9\,928$ & $64.64\%$\\
$100$ & $8$ & $60\,271$ & $26\,685$ & $55.72\%$\\
$1\,000$ & $10$ & $169\,157$ & $117\,011$ & $30.83\%$\\
$10\,000$ & $11$ & $1\,011\,962$ & $933\,686$ & $7.74\%$\\
\bottomrule
\end{tabular}
\caption{Sufficient row counts for $n=2^{30}$, $\alpha=\beta=1$,
$\varepsilon=1/2$, $\delta=1/100$, and
$q=\lceil\log_4(4d/\delta)\rceil$. These are certified prescriptions
valid for every fixed isometry $U$, not empirical thresholds or optimal
sample counts.}
\label{tab:srht-prefix-counts}
\end{table}

The displayed counts were located by a floating-point search, increased
by one percent, and then verified using rational arithmetic. Every square
root in a matrix entry was rounded upward to a rational with denominator
$10^{18}$; the square root used to lower-bound $\underline c$ was rounded
downward. The nonnegative vector recurrences were evaluated exactly.
For the ordered-prefix column, the certified upper bounds on
$d\underline c^{-2q}\|G^{[q]}\cdots G^{[1]}e_0\|_2^2/
(\delta\varepsilon^{2q})$ are, rounded upward,
$0.937586$, $0.918227$, $0.899509$, and $0.890415$.
The exact rational comparisons, rather than these decimal summaries,
certify sufficiency; the accompanying \texttt{checks/check\_srht.py} and
\texttt{checks/check\_srht\_terminal.py} reproduce the two columns.

The improvement does not change the local term
$\beta(d+8j^2)$, the two-round structure, or the confidence exponent of
the underlying theorem. No smaller universal scalar row coefficient is
claimed. The new contribution here is the explicit original-law
selector certificate and its ordered-prefix refinement; the row
contractions, positive-effects mechanism, and centered sampling coupling
are inherited from \cite{GradedSource,SRHTSource}.


\section{Khatri--Rao: covariance-sensitive boundaries and fourth-order calibration}
\label{sec:kr}

Fix an isometry $U:\mathbb R^r\to\mathbb R^{n_1}\otimes\mathbb R^{n_2}$.
For independent standard Gaussian vectors $g_i,h_i$, put
\[
 y_i=U^{\mathsf T}(g_i\otimes h_i),\qquad
 X_i=y_i y_i^{\mathsf T}-I_r,\qquad
 Z=\frac1m\sum_{i=1}^m X_i,\qquad W=\mathbb E X_1^2.
\]
Thus $\mathbb EX_i=0$, and $\|Z\|\le\varepsilon$ is the usual squared-norm
embedding guarantee on the fixed subspace $\operatorname{ran}U$. All traces
below are unnormalized. The coefficients and the subspace are fixed before
the sketch is drawn.

The original Khatri--Rao proof already propagates ordered, stage-dependent
comparison matrices and permits the exact value of $\|W\|$ in its activation
bounds~\cite{KRSource}. Our additional stationary comparison retains the entire
spectrum of $W$ at the vacuum boundary. A separate fourth-order calculation
retains the exact contribution of one original sample and improves the
Gaussian pairing contribution. Neither argument identifies the quartic
Gaussian polynomial $X_i$ with a linear Gaussian matrix series.

\subsection{An ordered first-return comparison}

We record the trace comparison used below, including its proof to make the
noncommutative step explicit. This is the covariance-spectrum comparison of
\cite[Theorem 4.3 and Corollary 4.4]{GradedSource}; the underlying renewal identity
is the classical block-resolvent expansion.

\begin{lemma}[Covariance at the boundary]\label{lem:kr-boundary}
Let $M$ be a finite-dimensional self-adjoint operator with decomposition
\[
 M=\begin{pmatrix}0&B^*\\B&C\end{pmatrix},\qquad V=B^*B,
\]
relative to a vacuum space $\mathcal H_0$ and its orthogonal complement.
Decompose the complement into sectors $\mathcal H_1,\ldots,\mathcal H_K$,
assume $\operatorname{ran}B\subseteq\mathcal H_1$, and suppose the symmetric
entrywise nonnegative matrix $G$ satisfies $\|P_iCP_j\|\le G_{ij}$.
For $v\ge0$, set
\[
 T(v)=\begin{pmatrix}0&\sqrt v\,e_1^{\mathsf T}\\
                       \sqrt v\,e_1&G\end{pmatrix}.
\]
For every integer $q\ge1$,
\begin{equation}\label{eq:kr-boundary-general}
 \operatorname{Tr}_{\mathcal H_0}(P_0M^{2q}P_0)
 \le\sum_{\lambda\in\operatorname{spec}V}(T(\lambda)^{2q})_{00}.
\end{equation}
Eigenvalues are counted with multiplicity. If $s\ge\|V\|$ and $s>0$, the
right side is at most $(\operatorname{Tr}V/s)(T(s)^{2q})_{00}$.
\end{lemma}

\begin{proof}
Put $K_j=P_0M^jP_0|_{\mathcal H_0}$ and $F_j=B^*C^{j-2}B$ for $j\ge2$.
Resolvent expansion, or splitting a vacuum walk at its first return, gives
\[
 K_0=I,\quad K_1=0,\qquad K_n=\sum_{j=2}^n F_jK_{n-j}.
\]
The order of the factors is retained. Write the polar decomposition
$B=QV^{1/2}$. Then
$F_j=V^{1/2}A_jV^{1/2}$, where
$A_j=Q^*C^{j-2}Q$ and $\|A_j\|\le c_j:=(G^{j-2})_{11}$.
The last inequality follows by block-norm propagation inside the complement;
$Q$ is a contraction with range in sector one.
For any composition $j_1+\cdots+j_\ell=n$ with $j_t\ge2$, cyclicity and
Schatten H\"older give
\[
 \left|\operatorname{Tr}(F_{j_1}\cdots F_{j_\ell})\right|
 =\left|\operatorname{Tr}(A_{j_1}V\cdots A_{j_\ell}V)\right|
 \le\left(\prod_{t=1}^{\ell}c_{j_t}\right)\operatorname{Tr}V^\ell.
\]
For the scalar matrix $T(v)$ the corresponding first-return weight is
$v c_j$. Expanding both renewals and summing the preceding estimate proves
\eqref{eq:kr-boundary-general}. In particular,
$f_q(v)=(T(v)^{2q})_{00}$ is a polynomial with nonnegative coefficients and
zero constant term. For $0\le v\le s$, $f_q(v)\le(v/s)f_q(s)$, proving the
last assertion. No Loewner substitution inside ordered products is used.
\end{proof}

\subsection{The original-law Khatri--Rao certificate}

For $q\ge1$, set $K=\min(q,m)$ and
\[
 \mathcal B_q(r)=e(27rq+192q^2+24q).
\]
Choose any certified $\nu\ge\|W\|$; the original Wick estimate permits
$\nu=4r+4$. For $v\ge0$, define a symmetric tridiagonal matrix
$T_{q,m}^{\nu}(v)$ on $0,\ldots,K$ by
\begin{align}
 (T_{q,m}^{\nu}(v))_{00}&=0,&
 (T_{q,m}^{\nu}(v))_{kk}&=\mathcal B_q(r)/m\quad(1\le k\le K),\notag\\
 (T_{q,m}^{\nu}(v))_{01}&=(T_{q,m}^{\nu}(v))_{10}=\sqrt v,&
 (T_{q,m}^{\nu}(v))_{k,k+1}
 &=\frac{\sqrt{(k+1)(m-k)\nu}}m\quad(1\le k<K).
 \label{eq:kr-T}
\end{align}
The symmetric counterparts of the last entries are understood. The positive
diagonal allowance accounts for repeated visits to an already active sample.

\begin{theorem}[Covariance-spectrum Khatri--Rao bound]\label{thm:kr-spectrum}
Let $\lambda_1,\ldots,\lambda_r$ be the eigenvalues of $W$. Then
\begin{align}
 \mathbb E\operatorname{Tr}Z^{2q}
 &\le\sum_{a=1}^r
   \left(T_{q,m}^{\nu}(\lambda_a/m)^{2q}\right)_{00}
 \label{eq:kr-spectrum}\\
 &\le\frac{\operatorname{Tr}W}{\nu}
   \left(T_{q,m}^{\nu}(\nu/m)^{2q}\right)_{00}.
 \label{eq:kr-effective}
\end{align}
Consequently, either right side being at most $\delta\varepsilon^{2q}$
is sufficient for $\mathbb P\{\|Z\|>\varepsilon\}\le\delta$.
\end{theorem}

\begin{proof}
Work in $\mathbb R^r\otimes\bigotimes_{i=1}^mL^2(g_i,h_i)$ and let
$\mathcal M_Z$ be multiplication by $Z$ on vector polynomials. Restrict to
functions even under each separate sign flip of $g_i$ and $h_i$. Compress
to total Hermite degree at most $4q$ and at most $q$ nonconstant sample
registers. Denote the projection by $\Pi_q$ and the compressed operator by
$M_q=\Pi_q\mathcal M_Z\Pi_q$.
This is a finite-dimensional symmetric matrix. Multiplication raises total
degree by at most four and activates at most one sample. Thus, for every
vacuum vector $v$ and $j\le q$,
$M_q^jv=\mathcal M_Z^jv$, and therefore
\[
 \mathbb E\operatorname{Tr}Z^{2q}
 =\operatorname{Tr}_{\mathbb R^r}(P_0M_q^{2q}P_0).
\]
This uses only finite polynomial moments; the full Gaussian multiplier need
not be bounded.

Grade by the number of nonconstant registers. Centering gives a zero vacuum
block. Distinct samples applied to vacuum are orthogonal, so its first-step
Gram is exactly $B^*B=W/m$, and $B$ lands in grade one. The fresh-sample
contraction of \cite[Chain C, ``Fresh-sample contraction'']{KRSource} gives
\[
 \|P_{k+1}M_qP_k\|
 \le\frac{\sqrt{(k+1)(m-k)\nu}}m.
\]
For completeness, each missing sample contributes a map with squared norm
at most $\nu/m^2$; an output set of size $k+1$ receives $k+1$ inputs, and
each input set has $m-k$ missing labels. Cauchy--Schwarz and orthogonality
give the displayed estimate on the full grade, not just on a basis vector.

For the diagonal block we import exactly the proved original-law bound
\cite[Chain C, ``Sum of already-active sample interactions'']{KRSource}:
$\|P_kM_qP_k\|\le\mathcal B_q(r)/m$ for $k\ge1$.
Its inputs are the local-degree estimate
\[
 \|y^{\mathsf T}f_d\|_2^2
 \le e\bigl[(2d+1)r+4d^2+2d\bigr]\|f_d\|_2^2
\]
and the fact that at most three input degrees overlap at any output degree.
Summing over the actual local degrees $d_i$, using
$\sum d_i\le4q$, $\sum d_i^2\le16q^2$, and at most $q$ active registers,
gives $\mathcal B_q(r)$. Centering contributes a lower bound $-qI$, also
covered by $\mathcal B_q(r)$. This retains the nonlinear repeated-sample
costs of the original proof. Apply Lemma~\ref{lem:kr-boundary} to the rest
matrix specified in \eqref{eq:kr-T}, then use trace Markov.
\end{proof}

If $W=sI_r$ and $\nu=s$, \eqref{eq:kr-spectrum} is exactly the old scalar
stationary certificate with exact variance. Thus flat covariance does not
illustrate an additional spectrum gain. A nonflat example does. Represent
$y_a=g^{\mathsf T}B_a h$ using
\[
 B_1=\operatorname{diag}(a,b),\qquad
 B_2=\frac1{\sqrt2}\begin{pmatrix}0&1\\1&0\end{pmatrix},\qquad
 a^2+b^2=1,\quad ab=\frac1{12}.
\]
These matrices are Frobenius-orthonormal. Wick's formula gives
$\mathbb Ey_1^4=107/12$, $\mathbb Ey_2^4=6$ and
$\mathbb Ey_1^2y_2^2=19/6$; the off-diagonal covariance vanishes by sign
symmetry. Hence
\[
 W=\operatorname{diag}(133/12,49/6),\qquad
 \|W\|=133/12,\qquad
 \frac{\operatorname{Tr}W}{\|W\|}=\frac{231}{133}<2.
\]
At $q=3$, $\varepsilon=1/2$, and $\delta=1/100$, the following are sufficient
row counts for the same stationary repeated-visit allowance:
\begin{center}
\begin{tabular}{lr}
\toprule
Activation/boundary information retained & Sufficient $m$\\
\midrule
Uniform allowance $4r+4$ at all edges & $10\,581$\\
Exact norm $\|W\|$ at all edges & $10\,403$\\
Exact spectrum at the boundary, exact norm inside & $10\,107$\\
\bottomrule
\end{tabular}
\end{center}
The accompanying rational verification checks each displayed integer and
failure of its predecessor. The certificate decreases for $m\ge2q$:
$1/m$ decreases, and each squared interior edge is proportional to
$(m-k)/m^2$, whose derivative is $(2k-m)/m^3\le0$ in this range.
These are comparisons of stationary certificates: the original manuscript
also has stronger ordered degree-allocation certificates, whose minimum
with \eqref{eq:kr-spectrum} remains a valid choice.

\subsection{Exact fourth-order calibration}

This next statement applies beyond Khatri--Rao. Let $X$ be a centered real
symmetric random matrix with finite fourth moment, let $X_1,\ldots,X_m$ be
independent copies, and put $Z=m^{-1}\sum_iX_i$. Choose a centered Gaussian
symmetric matrix $G=\sum_{\ell=1}^N\gamma_\ell A_\ell$ whose complete
matrix-entry covariance matches that of $X$. Such a representation exists
by factoring the covariance on the finite-dimensional real symmetric matrix
space. Then
\[
 W=\mathbb EX^2=\sum_\ell A_\ell^2,\qquad
 \kappa_G^2=\lambda_{\max}\bigl([A_jA_i]_{i,j=1}^N\bigr).
\]
The block matrix here is symmetric, even though its individual blocks need
not be. It is the crossed coefficient kernel appearing in the Gaussian
creation estimates of Bandeira et al.~\cite{Bandeira}.

\begin{proposition}[Original-component fourth-order certificate]
\label{prop:kr-fourth}
With $\mu_4=\mathbb E\operatorname{Tr}X^4$,
\begin{align}
 \mathbb E\operatorname{Tr}Z^4
 &=\frac{\mu_4+(m-1)\mathbb E\operatorname{Tr}G^4}{m^3}
 \label{eq:kr-fourth-exact}\\
 &\le\frac{\mu_4+(m-1)
    \bigl(2\operatorname{Tr}W^2+\kappa_G^2\operatorname{Tr}W\bigr)}{m^3}.
 \label{eq:kr-fourth-kappa}
\end{align}
\end{proposition}

\begin{proof}
Centering and independence eliminate words with a singleton sample label.
There are $m$ all-equal contributions. The words with two distinct labels
give $m(m-1)$ times the three covariance pairings, precisely the Gaussian
fourth moment. This proves \eqref{eq:kr-fourth-exact}, the trace specialization
of the fourth-order calibration in \cite[Lemma 8.2]{GradedSource}.
Wick expansion gives
\[
 \mathbb E\operatorname{Tr}G^4
 =2\operatorname{Tr}W^2+
   \sum_{i,j}\operatorname{Tr}(A_iA_jA_iA_j).
\]
Let $S$ be the block column with entries $A_i$ and let
$C=[A_jA_i]_{i,j}$. The last sum is $\operatorname{Tr}(S^*CS)$,
which is at most $\kappa_G^2\operatorname{Tr}(S^*S)
=\kappa_G^2\operatorname{Tr}W$. This proves the bound directly.
\end{proof}

The baseline Gaussian bound is $3\operatorname{Tr}W^2$; consequently
\eqref{eq:kr-fourth-kappa} improves it when
$\kappa_G^2\operatorname{Tr}W<\operatorname{Tr}W^2$.
The original component term $\mu_4$ remains exact. Matching the covariance of
the row $y$ instead of the centered matrix $yy^{\mathsf T}-I$ would not
justify this proposition.

\subsection{An analytic full-product benchmark}

Take the whole product space, $y=g\otimes h$ with
$g\sim N(0,I_a)$ and $h\sim N(0,I_b)$. Put
$r=ab$ and $L=(a+2)(b+2)$. Then
\begin{equation}\label{eq:kr-product-pairing}
 W=(L-1)I_r,\qquad
 \mathbb E\operatorname{Tr}G^4=r(2L^2+3L+3).
\end{equation}
To verify the second formula, let $X'$ be an independent copy of $X$.
Covariance matching and Wick expansion imply
\[
 \mathbb E\operatorname{Tr}G^4
 =3\operatorname{Tr}W^2-\tfrac12\mathbb E\|[X,X']\|_F^2.
\]
For rank-one matrices $P=yy^{\mathsf T}$ and $P'=y'y'^{\mathsf T}$,
\[
 \tfrac12\|[P,P']\|_F^2
 =\|y\|^2\|y'\|^2\langle y,y'\rangle^2-\langle y,y'\rangle^4.
\]
For independent $d$-dimensional standard Gaussian vectors $u,u'$, direct
conditioning yields
\[
 \mathbb E\|u\|^2\|u'\|^2\langle u,u'\rangle^2=d(d+2)^2,
 \qquad \mathbb E\langle u,u'\rangle^4=3d(d+2).
\]
The two factors separate, so the expected half-commutator square is
$rL(L-9)$. Substitution proves \eqref{eq:kr-product-pairing}.

We also give a small analytic certificate for the crossed kernel, so the
intermediate comparison does not depend on a floating-point eigenvalue.
Let $\{A_i^{(d)}\}$ be the GOE coefficient family
$\sqrt2 E_{jj}$ and $E_{jk}+E_{kj}$ for $j<k$. A Gaussian covariance match
for $X$ consists of the three coefficient families
\[
 A_i^{(a)}\otimes I_b,\qquad
 I_a\otimes A_j^{(b)},\qquad
 A_i^{(a)}\otimes A_j^{(b)}.
\]
Indeed, writing $gg^{\mathsf T}=I_a+A$ and $hh^{\mathsf T}=I_b+B$ gives
$X=A\otimes I_b+I_a\otimes B+A\otimes B$; the three summands are pairwise
uncorrelated, and $A,B$ have GOE covariance. Thus block-norm comparison gives
\begin{equation}\label{eq:kr-kappa-3x3}
 \kappa_G^2\le\lambda_{\max}
 \begin{pmatrix}
  2&\sqrt{(a+1)(b+1)}&2\sqrt{b+1}\\
  \sqrt{(a+1)(b+1)}&2&2\sqrt{a+1}\\
  2\sqrt{b+1}&2\sqrt{a+1}&4
 \end{pmatrix}.
\end{equation}
Here are the needed block norms. The GOE row map has squared norm $d+1$
because $\sum_i(A_i^{(d)})^2=(d+1)I_d$. Its crossed kernel, identified with
an operator on tensors $u_{ij,k}=u_{ji,k}$, sends
$u_{ij,k}$ to $u_{ik,j}+u_{jk,i}$, and hence has norm at most two by the
triangle inequality for coordinate permutations. The three diagonal blocks
therefore have norms at most $2,2,4$. The first off-diagonal block is a
product of two row/column maps, while each remaining block is the tensor
product of one crossed kernel and one row/column map. Their norms are the
off-diagonal allowances in \eqref{eq:kr-kappa-3x3}. Applying the scalar
Rayleigh bound to the norms of the three vector blocks proves the claim.

For $a=b=2$, $r=4$, $L=16$, and
\[
 W=15I_4,\qquad
 \kappa_G^2\le\frac{9+\sqrt{97}}2,\qquad
 \mathbb E\operatorname{Tr}G^4=2252.
\]
The kernel bound follows by separating the antisymmetric vector $(1,-1,0)$
from the remaining two-dimensional block of \eqref{eq:kr-kappa-3x3}.
Writing $t=\|y\|^2$, the rank-one identity gives
\[
 \operatorname{Tr}(yy^{\mathsf T}-I_4)^4=t^4-4t^3+6t^2-4t+4.
\]
Since $\mathbb Et^j=\prod_{k=0}^{j-1}(2+2k)^2$, this yields
$\mu_4=138612$. Therefore the exact fourth trace moment is
\begin{equation}\label{eq:kr-product-exact-fourth}
 \mathbb E\operatorname{Tr}Z^4
 =\frac{138612+2252(m-1)}{m^3}.
\end{equation}
At $\varepsilon=1/2$, $\delta=1/100$, fourth-order trace Markov gives:
\begin{center}
\begin{tabular}{lrr}
\toprule
Gaussian pairing allowance & Value & Sufficient $m$\\
\midrule
Generic $3\operatorname{Tr}W^2$ & $2700$ & $2104$\\
Crossed-kernel certificate & $1800+30(9+\sqrt{97})$ & $1974$\\
Exact covariance pairing & $2252$ & $1928$\\
\bottomrule
\end{tabular}
\end{center}
The attached verifier uses rational upper and lower bounds on $\sqrt{97}$
and checks both the stated count and its predecessor. These improvements
are respectively $6.18\%$ and $8.37\%$ relative to the generic fourth-order
certificate. They do not assert improvement over an optimized all-order
certificate, or a new universal constant in the Khatri--Rao embedding theorem.
The exact pairing formula is a direct covariance calculation; the comparison
with the crossed kernel records explicitly what the combined operator input
contributes.


\section{Independent spherical-product rows}
\label{sec:spherical-product}

The Gaussian Khatri--Rao guarantees extend to a different, bounded row
distribution by conditional expectation.  This extension uses independent
rows; the rows of a single Haar frame are not independent instances of this
model.

Fix positive integers $a,b,r$ and an isometry
$U\in\mathbb R^{ab\times r}$, $U^{\mathsf T}U=I_r$.  For each sample
$i=1,\ldots,m$, let $u_i\in S^{a-1}$ and $v_i\in S^{b-1}$ be independent
uniform unit vectors, independent also across samples.  Define
\[
 y_i^{\rm sph}=U^{\mathsf T}(\sqrt a\,u_i\otimes\sqrt b\,v_i),
 \qquad
 Z_{\rm sph}=\frac1m\sum_{i=1}^m
       \bigl(y_i^{\rm sph}(y_i^{\rm sph})^{\mathsf T}-I_r\bigr).
\]
Let $y_i^{\rm gau}=U^{\mathsf T}(g_i\otimes h_i)$ and $Z_{\rm gau}$
denote the corresponding rows and centered empirical covariance for
independent $g_i\sim N(0,I_a)$ and $h_i\sim N(0,I_b)$.  Both row models
are isotropic on $\mathbb R^r$.

\begin{theorem}[Convex transfer and exact spherical covariance]
\label{thm:spherical-transfer}
For every convex function $F$ on real symmetric $r\times r$ matrices
for which the expectations exist,
\begin{equation}\label{eq:spherical-convex}
 \mathbb E F(Z_{\rm sph})\le \mathbb E F(Z_{\rm gau}).
\end{equation}
In particular, for every integer $q\ge1$,
\begin{equation}\label{eq:spherical-trace-transfer}
 \mathbb E\operatorname{Tr}Z_{\rm sph}^{2q}
 \le \mathbb E\operatorname{Tr}Z_{\rm gau}^{2q}.
\end{equation}
For a single row, write
$W_{\star}=\mathbb E(y^{\star}(y^{\star})^{\mathsf T}-I_r)^2$.
Then, with $c_{a,b}=ab/((a+2)(b+2))$,
\begin{equation}\label{eq:spherical-covariance}
 W_{\rm sph}=c_{a,b}(W_{\rm gau}+I_r)-I_r.
\end{equation}
\end{theorem}

\begin{proof}
Gaussian polar decomposition gives a coupling
\[
 g_i=R_i u_i,\qquad h_i=S_i v_i,\qquad
 y_i^{\rm gau}=\sqrt{\tau_i}\,y_i^{\rm sph},\qquad
 \tau_i=\frac{R_i^2S_i^2}{ab},
\]
where $R_i^2\sim\chi_a^2$, $S_i^2\sim\chi_b^2$, and all radii are
independent of all directions.  Since $\mathbb E\tau_i=1$,
\[
 \mathbb E[Z_{\rm gau}\mid (u_i,v_i)_{i=1}^m]=Z_{\rm sph}.
\]
Conditional Jensen proves \eqref{eq:spherical-convex}.  For symmetric
$A$, $\operatorname{Tr}A^{2q}=\|A\|_{S_{2q}}^{2q}$ is convex,
which proves \eqref{eq:spherical-trace-transfer}.  Finally,
$\mathbb E\tau_i^2=(a+2)(b+2)/(ab)$ and isotropy give
\[
 W_{\rm gau}+I_r
 =\mathbb E\bigl(y^{\rm gau}(y^{\rm gau})^{\mathsf T}\bigr)^2
 =c_{a,b}^{-1}
   \mathbb E\bigl(y^{\rm sph}(y^{\rm sph})^{\mathsf T}\bigr)^2
 =c_{a,b}^{-1}(W_{\rm sph}+I_r).
\]
\end{proof}

\begin{corollary}[Transfer of moment-based embedding certificates]
\label{cor:spherical-ose}
If a Gaussian Khatri--Rao certificate proves
$\mathbb E\operatorname{Tr}Z_{\rm gau}^{2q}\le C_q(m)$ and
$C_q(m)\le\delta\varepsilon^{2q}$, then the spherical-product sketch
satisfies
\[
 \mathbb P(\|Z_{\rm sph}\|>\varepsilon)\le\delta.
\]
Consequently, every sufficient row count obtained by this moment method,
including the refinements above, remains valid for independent
spherical-product rows.
\end{corollary}

\begin{proof}
For a symmetric matrix $A$, $\|A\|^{2q}\le\operatorname{Tr}A^{2q}$.
Apply Markov's inequality and Theorem~\ref{thm:spherical-transfer}.
\end{proof}

The transfer concerns moment certificates, not a comparison of the two
failure probabilities at every threshold.  It introduces a new row model
without worsening those certificates.  Although
\eqref{eq:spherical-covariance} also gives $W_{\rm sph}\preceq W_{\rm gau}$,
substituting this smaller covariance into every Gaussian interior block
estimate would require a new proof: the multiplication laws have changed.
The argument uses classical radius--direction independence and Jensen's
inequality; the Gaussian certificates being transferred are those of
\cite{KRSource} and the present refinements.

\subsection{An exact fourth-moment certificate for the full product space}

For the full product space, fixed spherical radii permit a stronger
calculation than simply transferring the Gaussian certificate.

\begin{proposition}[Full-product spherical fourth moment]
\label{prop:spherical-fourth}
Let $r=ab$, $L=(a+2)(b+2)$, $U=I_r$, and
$x=yy^{\mathsf T}-I_r$ for one spherical-product row.  Then
\[
 W_{\rm sph}=(r-1)I_r,\qquad
 \beta_{a,b}:=\mathbb E\operatorname{Tr}x^4=(r-1)^4+r-1.
\]
Define
\[
 \gamma_{a,b}
 =3r(r-1)^2-r^3\left(1-\frac9L\right).
\]
For every positive integer $m$,
\begin{equation}\label{eq:spherical-exact-fourth}
 \mathbb E\operatorname{Tr}Z_{\rm sph}^4
 =\frac{\beta_{a,b}+(m-1)\gamma_{a,b}}{m^3}.
\end{equation}
\end{proposition}

\begin{proof}
Since $\|y\|^2=r$, the eigenvalues of $x$ are $r-1$ once and
$-1$ with multiplicity $r-1$.  This proves the value of $\beta_{a,b}$;
isotropy and $(yy^{\mathsf T})^2=r\,yy^{\mathsf T}$ prove the covariance.
Let $x'$ and $y'$ be independent copies.  Put
\[
 \Delta=\frac12\mathbb E\|[x,x']\|_F^2
 =\operatorname{Tr}W_{\rm sph}^2
       -\mathbb E\operatorname{Tr}(xx'xx').
\]
The rank-one commutator identity gives
\[
 \Delta=r^2\mathbb E\langle y,y'\rangle^2
              -\mathbb E\langle y,y'\rangle^4.
\]
For independent uniform unit vectors $u,u'\in S^{a-1}$,
\[
 \mathbb E\langle\sqrt a\,u,\sqrt a\,u'\rangle^2=a,
 \qquad
 \mathbb E\langle\sqrt a\,u,\sqrt a\,u'\rangle^4
       =\frac{3a^3}{a+2}.
\]
These standard spherical moments follow by conditioning on $u'$ and
using rotational invariance, or by Gaussian polar decomposition.
Multiplication of the two independent factor moments yields
$\mathbb E\langle y,y'\rangle^2=r$ and
$\mathbb E\langle y,y'\rangle^4=9r^3/L$.
Thus $\Delta=r^3(1-9/L)$.

In expanding the fourth trace moment of a sum of independent centered
$x_i$, only one-name words and two-name pairings survive.  Cyclicity
of trace gives
\[
 \mathbb E\operatorname{Tr}\left(\sum_{i=1}^m x_i\right)^4
 =m\beta_{a,b}
  +m(m-1)\left(2\operatorname{Tr}W_{\rm sph}^2
                     +\mathbb E\operatorname{Tr}(xx'xx')\right).
\]
The bracket is $3\operatorname{Tr}W_{\rm sph}^2-\Delta
=\gamma_{a,b}$.  Division by $m^4$ proves the result.
\end{proof}

For $a=b=2$, these quantities are $r=4$, $W_{\rm sph}=3I_4$,
$\beta_{2,2}=84$, and $\gamma_{2,2}=80$.  Therefore
\begin{equation}\label{eq:spherical-fourth-example}
 \mathbb E\operatorname{Tr}Z_{\rm sph}^4
       =\frac{80m+4}{m^3}.
\end{equation}
At $\varepsilon=1/2$ and $\delta=1/100$, the smallest integer satisfying
the resulting fourth-order Markov certificate is $m=358$:
\[
 1600(80\cdot358+4)\le358^3,
 \qquad
 1600(80\cdot357+4)>357^3.
\]
The certificate decreases with $m$.  The corresponding Gaussian
fourth-moment certificate in the preceding section requires $m=1928$.
This comparison quantifies the benefit of the spherical-product model
for this particular sufficient certificate; it is not a reduction of
the Gaussian theorem's constant or a comparison with the best available
all-order certificate.


\section{Reproducibility and use of the bounds}
The accompanying \texttt{checks} directory uses only the Python standard
library. The script \texttt{run\_checks.py} runs the SRHT ordered-prefix and
terminal comparisons, the Khatri--Rao covariance-spectrum recurrence,
the Gaussian fourth-order certificates, and the spherical fourth-order
calculation. Floating-point searches locate SRHT candidates; every reported
sufficiency claim is then checked by rational upper bounds. The
Khatri--Rao verifiers additionally check the preceding integer for the
specified monotone certificate. The spherical example is checked both
by exact integer threshold inequalities and by an independent explicit
covariance basis.

These numerical checks verify evaluation of the displayed certificates.
The probabilistic and operator comparisons are justified by the proofs
and the explicitly stated imported lemmas. For any fixed application,
one may take the minimum of independently valid sufficient moment bounds.
This applies, in particular, to the new stationary Khatri--Rao bound and
the earlier degree-allocated bounds, and to the transferred Gaussian
certificate and the direct spherical fourth-order certificate.

\appendix


\section{Conditional spherical linear pencils}
\label{sec:conditional-spherical}

For a linear matrix pencil, Gaussian polar decomposition gives an exact
moment correction rather than only the convex comparison of
Section~\ref{sec:spherical-product}.  This is a complementary extension:
the centered Khatri--Rao covariance error is nonlinear in its factors and
is not itself the linear pencil considered here.  Fresh spherical
directions in conditional frame arguments, such as those in
\cite{HSSSource}, provide one setting where the following hypotheses apply.

\begin{theorem}[Conditional spherical Jacobi comparison]
\label{thm:conditional-spherical-jacobi}
Condition on a sigma-field $\mathcal F$.  Let $u$ be conditionally uniform
on the real unit sphere $S^{h-1}$ and let
$B_1,\ldots,B_h$ be $\mathcal F$-measurable Hermitian $D\times D$ matrices.
Set
\[
 Y=\sqrt h\sum_{i=1}^h u_iB_i,\quad
 W=\sum_iB_i^2\ne0,\quad s=\|W\|,\quad
 \kappa^2=\lambda_{\max}([B_jB_i]_{i,j=1}^h).
\]
For $p\ge1$, let $J_p(v)$ be the symmetric tridiagonal matrix on
$0,\ldots,p$ with zero diagonal and edges
\[
 [J_p(v)]_{0,1}=\sqrt v,\qquad
 [J_p(v)]_{k,k+1}=\sqrt{s+k\kappa^2}\quad(1\le k<p).
\]
With $a_{h,p}=\prod_{j=0}^{p-1}h/(h+2j)$, almost surely,
\begin{align}
 \mathbb E[\operatorname{Tr}Y^{2p}\mid\mathcal F]
 &\le a_{h,p}\sum_{\lambda\in\operatorname{spec}(W)}
                    (J_p(\lambda)^{2p})_{00}\label{eq:spherical-jacobi}\\
 &\le a_{h,p}\frac{\operatorname{Tr}W}{s}
                    (J_p(s)^{2p})_{00}.\nonumber
\end{align}
Eigenvalues are counted with multiplicity.  At fourth order the sharper
bound
\begin{equation}\label{eq:spherical-linear-fourth}
 \mathbb E[\operatorname{Tr}Y^4\mid\mathcal F]
 \le\frac{h}{h+2}
          \left(2\operatorname{Tr}W^2+\kappa^2\operatorname{Tr}W\right)
\end{equation}
also holds.  If $W=0$, then $Y=0$ and the conclusions are interpreted
as zero without the displayed division by $s$.
\end{theorem}

\begin{proof}
Fix the conditioned coefficients and take independent standard real
Gaussians $g_i$.  Set $X=\sum_i g_iB_i$.  The decomposition
$g=Ru$, with $R$ independent of $u$ and $R^2\sim\chi_h^2$, gives the
matrix identity
\begin{equation}\label{eq:spherical-radial-identity}
 \mathbb E[Y^{2p}\mid\mathcal F]=a_{h,p}\mathbb E_gX^{2p}.
\end{equation}

For completeness, the interior Gaussian bound used in
\eqref{eq:spherical-jacobi} is the bosonic specialization of the creation
estimate in \cite[Proposition 4.1]{Bandeira}.  On Gaussian Hermite space
write $L=\sum_i B_i\otimes a_i^*$, where $a_i^*$ and $a_i$ are creation
and annihilation.  If $R_k f=(a_i f)_i$ on degree $k$, the canonical
commutation relation gives
\[
 L^*L=W\otimes I+R_k^*[B_jB_i]_{i,j}R_k,
 \qquad R_k^*R_k=kI.
\]
Consequently $\|L|_{\text{degree }k}\|^2\le s+k\kappa^2$.
The degree-$p$ cutoff of $L+L^*$ preserves its first $p$ vacuum iterates,
so its vacuum moment of order $2p$ equals $\mathbb E_gX^{2p}$.
The vacuum-to-degree-one Gram is exactly $W$.
The covariance-spectrum comparison of Lemma~\ref{lem:kr-boundary},
equivalently \cite[Corollary 4.4]{GradedSource}, therefore bounds this
Gaussian trace moment by the sum in \eqref{eq:spherical-jacobi}.
Its nonnegative return-polynomial relaxation gives the second inequality.
Multiplication by $a_{h,p}$ proves both spherical comparisons.

For the fourth-order refinement, Wick's formula gives
\[
 \mathbb E_g\operatorname{Tr}X^4
 =2\operatorname{Tr}W^2
       +\sum_{i,j}\operatorname{Tr}(B_iB_jB_iB_j).
\]
Regarding $(B_i)_i$ as a vector in the direct sum of Hilbert--Schmidt
spaces, its quadratic form against $[B_jB_i]_{i,j}$ is the last sum.
It is at most $\kappa^2\sum_i\operatorname{Tr}B_i^2
=\kappa^2\operatorname{Tr}W$.  Apply
\eqref{eq:spherical-radial-identity} at $p=2$.
\end{proof}

The fourth-order right side improves the $p=2$ Jacobi right side by
$\frac{h}{h+2}(s\operatorname{Tr}W-\operatorname{Tr}W^2)\ge0$.
The spherical factor is exact and equals one at second order.
Its role is to extend a Gaussian comparison to the stated conditional
law; it does not improve a proof already using exact conditional second
moments.  The ingredients have separate origins: classical polar
decomposition supplies the radial factor, \cite{Bandeira} supplies the
interior creation inequality, and \cite{GradedSource} supplies the
covariance-sensitive boundary comparison.  No joint-frame or accumulated
adaptive-process bound follows merely from the one-direction conditional
statement.


\begin{thebibliography}{9}
\addcontentsline{toc}{section}{References}
\bibitem{GradedSource}
D. Heidary.
\emph{Graded Multiplication Operators for Random-Matrix Moments},
version 11, 19 September 2026.
Theorem 4.3, Corollary 4.4, Lemmas 3.11--3.12, Lemma 8.2,
and the SRHT appendix.
\href{https://github.com/DiarHaidary/Results/blob/3a93d33c8b80326615b5b6e63654f3d84edf431d/Graded\%20Multiplication\%20Operators\%20for\%20Random-Matrix\%20Moments.pdf}{Public manuscript}.

\bibitem{SRHTSource}
D. Heidary.
\emph{Two-round SRHT: hybrid constants and proof-chain comparison}.
Retained row-factor, positive-effects, and fixed-size sampling arguments.
\href{https://github.com/DiarHaidary/Results/blob/3a93d33c8b80326615b5b6e63654f3d84edf431d/Wick-Hermite\%20Hybrid/srht_hybrid_constants_template.tex}{Public source}.

\bibitem{KRSource}
D. Heidary.
\emph{Sharp Gaussian Khatri--Rao Embeddings: Three Proof Chains:
Concentration, Framework, and Wick--Hermite}.
Consolidated proof-chain manuscript, Chain C: fresh-sample contraction,
sum of already-active sample interactions, and exact covariance examples.
\href{https://github.com/DiarHaidary/Results/blob/3a93d33c8b80326615b5b6e63654f3d84edf431d/Wick-Hermite\%20Hybrid/A\%20Tight\%20Analysis\%20of\%20Khatri-Rao\%20Oblivious\%20Subspace\%20Embeddings/khatri_rao_three_proof_chains_with_chainD_main.tex}{Public source}.

\bibitem{Bandeira}
A. S. Bandeira, D. Kunisky, P. Nizi\'{c}-Nikolac, L. Pesenti,
and R. Wang.
\emph{Matrix Concentration and Equivalent Operators on Fock Spaces}.
arXiv:2610.01982v1, 2026.
\href{https://arxiv.org/abs/2610.01982v1}{arXiv manuscript}.

\bibitem{HSSSource}
D. Heidary.
\emph{HSS: all routes and obstructions}.
Conditional frame laws and local Haar comparisons.
\href{https://github.com/DiarHaidary/Results/blob/3a93d33c8b80326615b5b6e63654f3d84edf431d/HSS/HSS\%20all\%20routes\%20and\%20obstructions.tex}{Public source}.
\end{thebibliography}
\end{document}
